\documentclass[11pt]{article}

\usepackage[margin=1in]{geometry}
\usepackage{amsmath,amssymb}
\usepackage{graphicx}
\usepackage{booktabs}
\usepackage{hyperref}
\usepackage{natbib}
\usepackage{xcolor}
\usepackage{caption}

\newcommand{\todo}[1]{\textcolor{red}{[TODO: #1]}}
\newcommand{\wq}{W^{Q}}
\newcommand{\wk}{W^{K}}
\newcommand{\wv}{W^{V}}

\title{Do Vision Transformers Need Three Projections?\\
Scaling Projection-Sharing Attention to ImageNet}

\author{
Ali Kayyam \\
BrainChip Inc., Laguna Hills, CA, USA \\
\texttt{akayyam@brainchip.com}
}

\date{Draft --- \today}

\begin{document}
\maketitle

\begin{abstract}
We study whether the projection-sharing attention variants shown to preserve
language-model quality while halving the KV cache
\citep{kayyam2026three} generalize to large-scale vision transformers.
Prior work in this line --- including our own systematic study of
query/key/value projection sharing on synthetic, vision (MNIST/CIFAR/TinyImageNet),
and language-modeling tasks, and a concurrent line of work reframing the
key-value-shared variant as ``Keyless Attention'' with a value-space routing
interpretation \citep{gao2026keyless} --- has validated projection sharing
at small-to-medium model and dataset scale. It remains open whether these gains persist at ImageNet scale, and whether
the overfitting-reduction effect attributed to gradient entanglement between
routing and retrieval parameters \citep{gao2026keyless} is a language-specific
artifact or a general property of shared-projection attention. We train
ViT-S/16 from scratch on full ImageNet-1k under a compressed DeiT-style
recipe, once each, for the QKV baseline and three projection-sharing
variants (Q-K=V, Q=K=V\(^+\), QVV(3)). Q-K=V and QVV(3) close most of the gap
to the QKV baseline (71.7\% and 72.5\% top-1 vs.\ 74.4\%), while the fully
symmetric Q=K=V\(^+\) variant lags substantially (66.8\%) --- a larger
degradation, in relative terms, than the 3.1\% perplexity cost reported for
the language-modeling setting, and one that does not support the hypothesis
that symmetric attention maps are less costly in vision than in causal
language modeling. We find no measurable overfitting-gap difference between
variants at this schedule: the heavily regularized 100-epoch DeiT recipe
(RandAugment, Mixup, CutMix, stochastic depth, label smoothing) leaves final-
and best-epoch validation loss within noise of each other for every
condition, including the QKV baseline, so the gradient-entanglement
regularization effect reported for language models is neither confirmed nor
refuted here. A follow-up run (Section~\ref{sec:followup}) closes a gap
flagged in the original design --- the Q=K-V\(^+\) condition, previously
skipped to fit four conditions across four GPUs --- and adds an unplanned
ablation removing the 2D positional correction from Q=K=V\(^+\). Q=K-V\(^+\)
reaches 70.4\% top-1, trailing its asymmetric analog Q-K=V by only 1.3 pts
(vs.\ Q=K=V\(^+\)'s 4.9-pt gap), suggesting the earlier symmetric-attention
penalty is driven more by collapsing \(V\) into the shared projection than
by attention-map symmetry itself; dropping the positional correction
entirely raises Q=K=V from 66.8\% to 70.7\%, though this run also doubled
the effective batch size, so the two effects are confounded and not yet
separable. A ViT-B/16 confirmatory run, an ImageNet-100 pilot phase, and a
downstream detection-transfer evaluation were part of the original design but
were descoped for this iteration; we report what a single-scale, single-seed
ViT-S/16 study on ImageNet-1k shows and leave the rest as future work.
\end{abstract}

\section{Introduction}
\label{sec:intro}

The query-key-value (QKV) formulation of self-attention \citep{vaswani2017attention}
is near-universal in both language and vision transformers, yet the individual
necessity of each of the three projections is not well understood. We
previously showed that in language models, merging the key and value
projections (\(K = V\), which we term \textbf{Q-K=V}), merging the query and
key projections with a corrective 2D positional term (\textbf{Q=K-V}\(^+\)),
and even collapsing all three into a single projection can match or nearly
match standard QKV attention, with the \(K{=}V\) variant achieving a 50\%
KV-cache reduction at only 3.1\% perplexity degradation on 300M--1.2B
parameter language models trained on 10B tokens \citep{kayyam2026three}.
That study also included a preliminary vision pass on MNIST, CIFAR, and
TinyImageNet.

A concurrent paper, \emph{Keyless Attention} \citep{gao2026keyless},
independently arrives at the same base architecture as our Q-K=V variant ---
algebraically, \(\mathrm{softmax}(QK^\top)K \equiv \mathrm{softmax}(QV^\top)V\)
when \(K=V\) --- but frames it as eliminating the key projection and routing
directly through value space, rather than as a special case of projection
sharing. That paper extends the idea along a new axis, \emph{depth-\(m\)
attention factorization}, in which the query and/or value projection is
itself a composition of \(m\) learned matrices, and shows that the depth-3
instantiation (\textbf{QVV(3)}) matches the parameter count of standard QKV
attention while still requiring only a value cache at inference time. It
also reports a consistent finding across five language models and four
architectures (GPT-2, Pythia, Qwen2, Llama 3.2): projection-shared attention
degrades \emph{more slowly} after the best-validation epoch than standard
attention, and attributes this to a gradient-entanglement effect between the
routing and retrieval parameters that acts as an implicit regularizer. All
five of their experiments are causal language models; QVV(2) and QVV(3) are
never evaluated on a vision transformer, so the equivalence to our Q-K=V
variant is algebraic rather than empirically confirmed on non-causal,
2D-grid attention --- the gap this paper's ImageNet-1k experiments address.

Two questions follow naturally from combining these two lines of work, and
motivate this paper:

\begin{enumerate}
    \item \textbf{Does projection sharing hold at ImageNet scale?} All
    vision evidence to date for this family of methods is at
    MNIST/CIFAR/TinyImageNet scale, with small models. Vision transformers
    at the scale used in production (ViT-B/16, ViT-L/16, ImageNet-1k) have
    a very different data/parameter regime, and it is not obvious that a
    5.6\%-perplexity-scale finding on WikiText-103 will transfer.
    \item \textbf{Is the overfitting-reduction effect language-specific?}
    The gradient-entanglement argument in \citet{gao2026keyless} is
    architecture-agnostic in principle, but has only been tested on
    autoregressive causal language modeling. Vision transformers use
    bidirectional (non-causal) attention over a 2D patch grid, which
    changes both the routing geometry and the overfitting dynamics
    (ImageNet-1k, at 1.28M images, is also a substantially larger and more
    diverse training set relative to model capacity than WikiText-103 at
    equivalent scale).
\end{enumerate}

We additionally note that vision transformers offer a structural advantage
absent in language: patches sit on an explicit 2D grid, so the 2D
positional-encoding correction we introduced for the symmetric Q=K-V and
Q=K=V variants in \citet{kayyam2026three} is a \emph{native} fit rather than
an artificial addition, unlike in the 1D sequential setting of language
modeling.

This paper originated as a design document laying out architectures,
training recipes, compute budget, hypotheses, and evaluation protocol for a
multi-scale study (Section~\ref{sec:setup} still describes that original
plan). For this iteration, the scope was deliberately narrowed to a single
architecture (ViT-S/16), a single training run per condition, directly on
full ImageNet-1k with no ImageNet-100 pilot phase, and no detection-transfer
experiment --- prioritizing a complete, correctly-run single-scale result
over a partially-completed multi-scale one. Section~\ref{sec:results} reports
the resulting numbers; Section~\ref{sec:limitations} discusses what the
descoping leaves untested.

\subsection{Contributions}
\begin{itemize}
    \item Evaluation of three projection-sharing attention variants (Q-K=V,
    Q=K=V\(^+\), and QVV(3)) against a standard QKV baseline on ViT-S/16
    trained from scratch on full ImageNet-1k, extending prior evidence for
    this family of methods beyond MNIST/CIFAR/TinyImageNet scale.
    \item A direct test of whether the overfitting-reduction /
    gradient-entanglement effect reported for language models
    \citep{gao2026keyless} reproduces under non-causal, 2D-grid attention ---
    which, at this training schedule, it does not measurably do, for any
    variant including the baseline (Section~\ref{sec:discussion}).
    \item A follow-up run adding the previously-skipped Q=K-V\(^+\)
    condition and an unplanned ablation of the 2D positional correction on
    Q=K=V (Section~\ref{sec:followup}), refining the H3 verdict from
    Section~\ref{sec:discussion}.
    \item \emph{Deferred to future work:} a ViT-B/16 (and ViT-L/16)
    confirmatory run at larger scale, an ImageNet-100 pilot phase, and a
    downstream detection-transfer evaluation (FCOS/GFL-style), all part of
    the original design (Section~\ref{sec:setup}) but out of scope here.
\end{itemize}

\section{Related Work}
\label{sec:related}

\paragraph{Projection sharing in attention.}
\citet{kayyam2026three} systematically evaluate three projection-sharing
constraints --- Q=K-V, Q-K=V, and Q=K=V --- and show that the last two,
which produce symmetric attention maps, can be de-symmetrized via a 2D
sinusoidal positional correction. The Q-K=V variant is highlighted as the
most favorable trade-off, achieving 50\% KV-cache reduction with a 3.1\%
perplexity cost. \citet{gao2026keyless} re-derive the same base mechanism
from a ``value-space routing'' motivation and generalize it via depth-\(m\)
attention factorization, introducing the QVV(3) variant that matches
standard attention's parameter count and, in their experiments, matches or
exceeds it on perplexity and zero-shot benchmarks across five models while
halving the KV cache. \citet{borji2023kv} is an earlier, independent
precursor examining a KV-only (query-key merged) transformer formulation.

\paragraph{Shared projections in vision transformers.}
\citet{lee2021rethinking} share nonlinear embedding layers across Q, K, V
in tiny ImageNet-scale ViTs (3.1M parameters) with a code-parameterized
sharing scheme, reporting 71.4\% top-1 accuracy at small scale. MASA
\citep{masa2025} decomposes attention projection matrices into a shared
dictionary of matrix atoms, cutting attention parameters by 66.7\% with
on-par performance, and includes a ViT classification experiment. Neither
targets KV-cache reduction or ViT-B/L scale training on full ImageNet-1k;
both motivate parameter reduction rather than routing/retrieval unification,
and neither reports an overfitting/regularization analysis. We view this as
the primary novelty gap this study addresses on the vision side.

\paragraph{KV-cache reduction more broadly.}
Multi-Query Attention \citep{shazeer2019mqa} and Grouped-Query Attention
\citep{ainslie2023gqa} reduce the KV cache by sharing keys/values across
attention heads, and Multi-Head Latent Attention \citep{deepseek2024mla}
compresses keys and values into a shared low-rank latent. These methods are
orthogonal to projection sharing, which instead removes redundancy between
the key and value \emph{projections themselves}; \citet{kayyam2026three} and
\citet{gao2026keyless} both note the two families compose.

\paragraph{ViT training recipes.}
We adopt training conventions from \citet{dosovitskiy2021vit} and
\citet{touvron2021deit}; DeiT's data-efficient training recipe
(RandAugment, Mixup, CutMix, stochastic depth, label smoothing) is
essential for training ViT-B from scratch on ImageNet-1k without
ImageNet-21k pretraining, and we use it as our default recipe with a
compressed epoch budget (Section~\ref{sec:setup}).

\section{Method}
\label{sec:method}

\subsection{Recap: projection-sharing variants}
Standard multi-head attention computes, per head,
\begin{equation}
    \mathrm{Attn}(Q,K,V) = \mathrm{softmax}\!\left(\frac{QK^\top}{\sqrt{d_k}}\right)V,
    \qquad Q=X\wq,\ K=X\wk,\ V=X\wv .
\end{equation}
We evaluate the following variants, applied to a standard ViT patch-embedding
backbone (Section~\ref{sec:setup}):

\begin{description}
    \item[Q-K=V.] Set \(V=K\), i.e.\ drop \(\wv\) and read the output from
    the key projection: \(A=\mathrm{softmax}(\alpha QK^\top)K\). Equivalent
    to the QVV(2) base case of \citet{gao2026keyless}. Halves the
    cache (single projection cached in the autoregressive setting; for
    ViT, halves activation memory for the shared K/V tensor during
    training and any cached-inference use case, e.g.\ streaming or
    tiled inference).
    \item[Q=K-V\(^+\).] Set \(Q=K\), producing a symmetric attention map
    \(\mathrm{softmax}(\alpha KK^\top)V\); corrected for asymmetry with a
    2D positional term (Section~\ref{sec:posenc}).
    \item[Q=K=V\(^+\).] Collapse all three projections to one, with the
    same 2D positional correction.
    \item[QVV(3).] Depth-3 factorization from \citet{gao2026keyless}:
    \(Q = X W^{Q_1}W^{Q_2}\), \(V=X\wv\), output
    \(\mathrm{softmax}(QV^\top/\sqrt{d_k})V\). Matches standard attention's
    parameter count; at inference, \(W^{Q_1}W^{Q_2}\) is fused into a single
    matrix, so it incurs no extra cost over Q-K=V while retaining full
    QKV parameter capacity.
\end{description}

All variants preserve the standard ViT patch embedding, learned or
sinusoidal position embedding for the base sequence order, MLP block,
LayerNorm placement, and classification head; the attention block is the
sole controlled variable, matching the ablation design of both source
papers.

\subsection{2D positional correction on a native 2D grid}
\label{sec:posenc}
\citet{kayyam2026three} introduce a 2D sinusoidal positional correction to
de-symmetrize the Q=K-V and Q=K=V attention maps in language modeling, where
the 2D structure is an artificial construction imposed on a 1D token
sequence. In ViT, patches are natively arranged on a 2D grid
(\(H/p \times W/p\) for patch size \(p\)), so the same correction can be
applied using the model's actual patch coordinates rather than a synthetic
2D reindexing of a 1D sequence:
\begin{equation}
    \tilde{S}_{ij} = S_{ij} + \mathrm{Conv}_{1\times1}\big(P(\mathrm{row}_i,\mathrm{col}_i,\mathrm{row}_j,\mathrm{col}_j)\big),
\end{equation}
where \(P\) is a fixed 2D relative-position sinusoidal embedding indexed by
true patch row/column, and \(S\) is the symmetric raw attention score matrix.
We view this as a stronger test of the correction's design intent than the
language setting, and a natural point of comparison to existing 2D relative
position bias schemes in windowed vision transformers (e.g.\ Swin
\citep{liu2021swin}), which we will cite as a baseline design pattern rather
than a competing method.

\subsection{Hypotheses}
\label{sec:hypotheses}
\begin{description}
    \item[H1 (Accuracy parity).] Q-K=V and QVV(3) match standard ViT-B/16
    top-1 accuracy on ImageNet-1k within noise (\(\lesssim\)0.5 pt),
    consistent with the language-modeling degradation being small
    (3.1\%) and vision transformers being comparatively more
    over-parameterized relative to ImageNet-1k than LLMs are relative to
    their pretraining corpora.
    \item[H2 (Overfitting reduction transfers).] Shared-projection variants
    show a smaller train/val generalization gap than standard QKV at
    matched epoch count, replicating the gradient-entanglement
    regularization effect from \citet{gao2026keyless} outside the causal
    LM setting.
    \item[H3 (Symmetric variants underperform more in vision than language).]
    Because 2D spatial relationships in images are less inherently
    directional than causal token order in language, we expect the gap
    between corrected symmetric variants (Q=K-V\(^+\), Q=K=V\(^+\)) and
    the asymmetric variants (Q-K=V, QVV(3)) to be \emph{smaller} in vision
    than was reported in \citet{kayyam2026three} for language --- i.e.\
    symmetry is less costly for spatial patch attention than for causal
    sequence attention. This is falsifiable and could go the other way if
    fine-grained spatial ordering (e.g.\ for texture or object boundary
    localization) turns out to matter more than expected.
    \item[H4 (Transfer to detection is architecture-sensitive).] We expect
    accuracy parity on classification to \emph{not} guarantee parity on
    detection transfer, since detection localization loss is more
    sensitive to fine spatial attention structure; this is the primary
    motivation for including the downstream transfer experiment
    (Section~\ref{sec:transfer}) rather than treating ImageNet top-1 as
    sufficient evidence.
\end{description}

\section{Experimental Setup}
\label{sec:setup}

\subsection{Architectures}
\begin{itemize}
    \item \textbf{ViT-S/16} (22M params, 12 layers, hidden 384, 6 heads):
    primary model for the variant sweep, chosen for fast iteration.
    \item \textbf{ViT-B/16} (86M params, 12 layers, hidden 768, 12 heads):
    confirmatory run on the winning variant(s) from the sweep, and the
    scale most directly comparable to production vision-backbone usage.
    \item \textbf{ViT-L/16} (307M params): stretch goal, only if the
    ViT-S~$\to$~ViT-B trend is informative and compute allows; used to
    test whether any observed effect (H1--H3) grows or shrinks with
    scale, mirroring the depth-scaling analysis in
    \citet{kayyam2026three} (12-layer vs.\ 36-layer GPT-2).
\end{itemize}

\subsection{Datasets and schedule}
Two-phase plan to manage compute cost on 1--2 A100-class GPUs:
\begin{enumerate}
    \item \textbf{Pilot sweep --- ImageNet-100} (100-class subset,
    $\sim$130K images): all five conditions (QKV baseline, Q-K=V,
    Q=K-V\(^+\), Q=K=V\(^+\), QVV(3)) on ViT-S/16, 100 epochs, DeiT
    recipe. Used to pick the variant(s) that proceed to the confirmatory
    run and to get an early read on H2/H3.
    \item \textbf{Confirmatory run --- ImageNet-1k} (1.28M images): QKV
    baseline vs.\ the 1--2 best pilot variants, on ViT-B/16 (and ViT-L/16
    if feasible), 100--150 epoch compressed DeiT recipe (RandAugment,
    Mixup~$\alpha{=}0.8$, CutMix~$\alpha{=}1.0$, stochastic depth,
    label smoothing 0.1), rather than the full 300-epoch recipe, following
    the same ``controlled comparison over SOTA-chasing'' framing used for
    the WikiText-103 subset in \citet{kayyam2026three}.
\end{enumerate}

\begin{table}[h]
\centering
\caption{Estimated wall-clock training budget (bf16, DeiT-style recipe,
data pipeline assumed non-bottlenecked via FFCV/DALI). Two-GPU figures
assume $\sim$1.8$\times$ scaling from single-GPU due to gradient-sync
overhead.}
\label{tab:compute}
\begin{tabular}{@{}llrrr@{}}
\toprule
Dataset & Model & Epochs & 1$\times$ A100 & 2$\times$ A100 \\
\midrule
ImageNet-100 & ViT-S/16 & 100 & 5--7 hrs & $\sim$3 hrs \\
ImageNet-100 & ViT-B/16 & 100 & 15--20 hrs & 8--10 hrs \\
ImageNet-1k  & ViT-S/16 & 100--150 & 12--18 hrs & 6--9 hrs \\
ImageNet-1k  & ViT-B/16 & 100--150 & 2--3 days & 1--1.5 days \\
ImageNet-1k  & ViT-B/16 & 300 (full DeiT) & 6--8 days & 3--4.5 days \\
\bottomrule
\end{tabular}
\end{table}

With 5 conditions in the ImageNet-100 pilot and 2--3 in the ImageNet-1k
confirmatory phase, total estimated compute is on the order of
\(\sim\)4--5 GPU-days on a single A100, or \(\sim\)2.5--3 GPU-days on two,
dominated by the ImageNet-1k ViT-B/16 runs.

\paragraph{What was actually run.} The plan above was descoped before
execution: we ran only the ImageNet-1k phase, only on ViT-S/16, and only one
seed per condition, skipping the ImageNet-100 pilot and the ViT-B/16/ViT-L/16
confirmatory scale-up. Of the five conditions in the design, four were run
--- QKV baseline, Q-K=V, Q=K=V\(^+\), and QVV(3) --- omitting Q=K-V\(^+\)
(query-key merged with positional correction) to keep the sweep to what
fits on four GPUs at once. Each condition trained for 100 epochs on the
compressed DeiT recipe of Section~\ref{sec:setup} (RandAugment, Mixup
\(\alpha{=}0.8\), CutMix \(\alpha{=}1.0\), stochastic depth, label smoothing
0.1), batch size 128 (the largest that fit in 11GB), on hardware substantially
below the A100-class budget in Table~\ref{tab:compute}: four RTX 2080 Ti GPUs,
one condition per GPU, run to completion in parallel. Measured throughput was $\sim$575 img/s/GPU, $\sim$3320s/epoch, i.e.\
$\sim$92 GPU-hours per condition; running all four conditions concurrently,
one per GPU, kept total wall-clock to the same $\sim$92 hours ($\sim$3.8
days) rather than requiring $4\times$ that.

\paragraph{Follow-up run.} A second batch, reported in
Section~\ref{sec:followup}, trained the two remaining conditions: Q=K-V\(^+\)
(omitted above) and an unplanned ablation, Q=K=V with the 2D positional
correction removed. With only two conditions to run and all four GPUs free,
each was trained 2-GPU DDP (batch size 128 per GPU, effective global batch
256, base LR linearly scaled to match) rather than the single-GPU,
batch-128 setup used for the other four conditions --- roughly
$2\times$ the effective batch size and LR of Table~\ref{tab:imagenet1k}, with
epoch time falling to $\sim$1630--1645s accordingly. This is a genuine
setting difference, not just a resource-allocation detail, so
Section~\ref{sec:followup}'s numbers are not a controlled addition to
Table~\ref{tab:imagenet1k}; we flag comparisons against it explicitly where
the batch-size difference matters.

\subsection{Evaluation metrics}
\begin{itemize}
    \item Top-1 / top-5 accuracy at best-validation epoch (mirrors
    Table~1 of \citet{kayyam2026three} and Table~1 of
    \citet{gao2026keyless}).
    \item Overfitting gap \(\Delta = \) (final-epoch val.\ loss) $-$
    (best-epoch val.\ loss), directly testing H2.
    \item Activation/cache memory reduction for the shared-projection
    variants, reported both in absolute terms and as a fraction of
    standard-attention memory (target: 50\% for Q-K=V and QVV(3), as in
    prior work).
    \item Throughput (images/sec, batch size fixed) at inference,
    matching the decode-throughput benchmarking methodology of
    \citet{gao2026keyless} adapted to non-autoregressive ViT inference.
\end{itemize}

\subsection{Downstream transfer: detection}
\label{sec:transfer}
To test H4, we fine-tune the ImageNet-1k-pretrained ViT-B/16 backbones
(QKV baseline and best-performing shared variant) into a detection
pipeline combining FCOS-style dense prediction, GFL loss, and SimOTA label
assignment, consistent with our existing detection scaffold. We report
COCO-style AP, AP\textsubscript{50}, AP\textsubscript{75}, and
size-stratified AP (small/medium/large), since we expect any localization
sensitivity from projection sharing to appear disproportionately on
small-object AP, where fine spatial routing matters most.

\section{Results}
\label{sec:results}

\subsection{ImageNet-1k, ViT-S/16}
Table~\ref{tab:imagenet1k} reports top-1/top-5 accuracy at the best
validation epoch, the overfitting gap \(\Delta\) (final-epoch minus
best-epoch validation loss), inference-time KV/activation cache relative to
the QKV baseline, parameter count, and per-GPU training batch size, for six
conditions trained to 100 epochs on full ImageNet-1k. Best-epoch numbers
(not final-epoch) are used for Top-1/Top-5 since a small amount of post-peak
drift is visible for every condition; final-epoch numbers, used only for
\(\Delta\), differ from the best-epoch numbers by at most 0.03 pts top-1 for
any condition. The first four rows are the original controlled sweep, one
condition per GPU at batch 128; the last two (below the rule) are a
follow-up batch trained together at effective batch 256 (2-GPU DDP,
linearly-scaled LR) rather than batch 128, discussed separately in
Section~\ref{sec:followup} because of that setting difference.

\begin{table}[h]
\centering
\caption{ImageNet-1k, ViT-S/16, 100 epochs, one run per condition. Top-1/
Top-5 at best validation epoch; \(\Delta\) Overfit is final-epoch minus
best-epoch validation loss (negative values indicate validation loss was
still falling at epoch 100). The last two rows trained at a different
effective batch size (Section~\ref{sec:followup}) and are not directly
controlled against the first four.}
\label{tab:imagenet1k}
\begin{tabular}{@{}lcccccc@{}}
\toprule
Variant & Top-1 & Top-5 & $\Delta$ Overfit & Cache & Params & Batch \\
\midrule
QKV (baseline) & 74.38\% & 91.93\% & $+$0.0010 & 100\% & 22.05M & 128 \\
Q-K=V          & 71.67\% & 90.32\% & $-$0.0004 & 50\%  & 20.28M & 128 \\
Q=K=V\(^+\)    & 66.76\% & 87.14\% & $-$0.0014 & 33\%  & 18.50M & 128 \\
QVV(3)         & 72.55\% & 90.81\% & $+$0.0009 & 50\%  & 22.05M & 128 \\
\midrule
Q=K-V\(^+\)     & 70.40\% & 89.47\% & $-$0.0011 & 50\% & 20.28M & 256 \\
Q=K=V (no $^+$) & 70.74\% & 89.74\% & $-$0.0027 & 33\% & 18.50M & 256 \\
\bottomrule
\end{tabular}
\end{table}

Relative to the QKV baseline, Q-K=V and QVV(3) close most but not all of the
gap (\(-\)2.71 and \(-\)1.83 pts top-1), while Q=K=V\(^+\) lags well behind
(\(-\)7.62 pts). Every \(\Delta\)Overfit value is within 0.0014 nats of zero
--- an order of magnitude smaller than would be needed to distinguish
variants, and consistent with none of the six models overfitting in any
meaningful sense at this schedule (see Section~\ref{sec:discussion}).
Cache and parameter figures match the design targets in
Section~\ref{sec:setup}: QVV(3) matches QKV's parameter count while
retaining Q-K=V's 50\% inference cache, and Q=K=V\(^+\) and Q=K=V (no
\(^+\)) tie for the fewest parameters of any condition tested by dropping
two of three projections.

\subsection{Follow-up: Q=K-V\(^+\) and a positional-correction ablation}
\label{sec:followup}
The last two rows of Table~\ref{tab:imagenet1k} are the two conditions
trained in the follow-up batch described in Section~\ref{sec:setup}:
Q=K-V\(^+\), the condition omitted from the original four-way sweep, and
Q=K=V (no \(^+\)), an ablation that removes the 2D positional correction
from Q=K=V\(^+\) while leaving every other projection unchanged. Both
trained 2-GPU DDP at effective batch 256 (vs.\ batch 128 for the first four
rows), so absolute numbers here are not directly controlled against them.

\paragraph{Closing the H3 gap.} Q=K-V\(^+\) and Q-K=V have identical
parameter counts and cache targets (Table~\ref{tab:imagenet1k}), differing
only in whether \(Q\) is a separate projection (Q-K=V) or tied to \(K\)
(Q=K-V\(^+\), corrected for the resulting symmetric attention map). Their
gap is 1.27 pts (71.67\% vs.\ 70.40\%), far smaller than Q=K=V\(^+\)'s 4.91-pt
gap to the same Q-K=V reference. Since Q=K=V\(^+\) differs from Q=K-V\(^+\)
by also collapsing \(V\) into the shared projection, this points at
\emph{that} collapse --- not the symmetric attention map itself --- as the
main source of the penalty reported in Section~\ref{sec:discussion}, more
consistent with H3's original motivation than the all-symmetric-variants
verdict the single Q=K=V\(^+\) data point supported. The batch-size caveat
above applies: Q=K-V\(^+\) and Q-K=V were not trained under matched
settings, so this reading should be treated as suggestive pending a
same-batch confirmatory run.

\paragraph{Positional correction ablation.} Removing the 2D positional
correction raises Q=K=V's top-1 from 66.76\% (Q=K=V\(^+\), Table~\ref{tab:imagenet1k})
to 70.74\% --- the opposite of what the correction was designed to do. This
is the largest single swing in Table~\ref{tab:imagenet1k}, but it is fully confounded with
the batch-256/higher-LR setting (Section~\ref{sec:setup}): we cannot
currently attribute it to the correction, the batch size, the LR, or an
interaction among them. We report it because of its size, not because we
can explain it; isolating the cause needs a same-batch-size run of Q=K=V\(^+\)
(or, symmetrically, Q=K=V at batch 128), left for future work.

\subsection{Post-hoc K=V surgery and distillation recovery}
\label{sec:surgery}
Table~\ref{tab:imagenet1k} asks whether \emph{jointly training} Q-K=V from a
random initialization matches QKV. A companion question, developed in a
parallel note on post-hoc weight surgery (same repository, \texttt{distillation}
branch), is different: can an already-trained QKV teacher's independent
$K$/$V$ projections be tied together \emph{after the fact}, with no
retraining, and if that collapses (it does, everywhere it has been tested so
far --- synthetic sequence tasks, MNIST/FMNIST/CIFAR-10/100, and a
300M-parameter FineWeb-Edu language model), how much of the gap does a short
distillation fine-tune against the QKV teacher's soft logits recover,
relative to the jointly-trained Q-K=V ceiling in Table~\ref{tab:imagenet1k}?
We repeat that protocol here as a first ImageNet-1k/ViT-S/16 data point,
reusing the QKV and Q-K=V checkpoints already trained for
Table~\ref{tab:imagenet1k} rather than training anything new: three surgery
modes (\textsc{keep\_k}: student's shared projection $\leftarrow$ teacher's
$W^K$; \textsc{keep\_v}: $\leftarrow W^V$; \textsc{avg}: $\leftarrow
\tfrac12(W^K+W^V)$), each evaluated zero-shot immediately after surgery, then
fine-tuned for 5 epochs against the QKV teacher's soft logits (cross-entropy
on hard labels + KL divergence at $T{=}2$, equal weighting), with light
augmentation (RandomResizedCrop + flip only --- the heavy DeiT recipe of
Section~\ref{sec:setup} is for training 100 epochs from scratch, not a short
recovery fine-tune, and Mixup/CutMix's soft labels do not compose with
hard-label CE).

\begin{table}[h]
\centering
\caption{Post-hoc K=V surgery on the ImageNet-1k QKV teacher (74.38\% top-1),
zero-shot and after 5 epochs of distillation, against the jointly-trained
Q-K=V ceiling (71.67\% top-1) from Table~\ref{tab:imagenet1k}.}
\label{tab:surgery}
\begin{tabular}{@{}lcccccc@{}}
\toprule
Mode & ZS Top-1 & ZS Top-5 & Distilled Top-1 & Distilled Top-5 & $\Delta$ vs.\ ceiling \\
\midrule
\textsc{keep\_k} & 0.08\% & 0.46\% & 71.01\% & 89.98\% & $-$0.66 pt \\
\textsc{keep\_v} & 1.91\% & 6.37\% & 71.93\% & 90.43\% & $+$0.26 pt \\
\textsc{avg}     & 0.66\% & 2.50\% & 71.87\% & 90.42\% & $+$0.20 pt \\
\bottomrule
\end{tabular}
\end{table}

All three collapse to near-random ($1/1000 = 0.1\%$ chance level) under
zero-shot surgery, consistent with every prior setting and with the
convex-hull argument motivating that note: forcing $K=V$ post-hoc confines
each attention block's output to the convex hull of vectors an
already-specialized $K$-projection was never trained to preserve query-independent
content in. The zero-shot ordering does not replicate the small-scale/LLM
pattern where \textsc{avg} was uniformly safest: here \textsc{keep\_v}
(1.91\% top-1) is safest, with \textsc{avg} (0.66\%) between it and
\textsc{keep\_k} (0.08\%) rather than above both.

After 5 epochs of distillation, though, the picture that matters more ---
recovery relative to the from-scratch ceiling --- looks \emph{better} here
than at any prior scale, not worse: all three modes land within 0.7 pt of
the 71.67\% ceiling, and \textsc{keep\_v} and \textsc{avg} both slightly
\emph{exceed} it (by 0.26 and 0.20 pt respectively), echoing the one
exception already on record at small scale (\textsc{avg}/CIFAR-10 also
edged past its from-scratch ceiling). Most notably, \textsc{keep\_k} ---
the mode with a \emph{persistent}, only-half-closed gap even after 10
epochs on CIFAR-10/100, and the worst-recovering mode at LLM scale too ---
reaches 71.01\% here, just 0.66 pt behind, with per-epoch gains
(+51.1, +7.3, +5.6, +4.6, +1.3 pts) decelerating sharply enough by epoch 5
that we did not judge a 10-epoch extension necessary. Both the zero-shot
severity and the post-distillation \textsc{keep\_k} penalty that motivated
treating it as the qualitatively harder recovery case at every prior scale
are much smaller at ImageNet-1k/ViT-S/16 scale --- an open question for why
(more distillation steps per epoch on a 1.28M-image training set than on
CIFAR's 50k, a wider/deeper model with more redundant capacity to route
around a mismatched basis, or something scale-specific to vision
classification) that this single run cannot settle on its own.

\subsection{ViT-B/16 confirmatory run}
Not run in this iteration; deferred to future work along with the
ImageNet-100 pilot phase (Section~\ref{sec:setup}).

\subsection{Detection transfer}
Not run in this iteration. No FCOS/GFL detection scaffold currently exists
in this codebase, and standing one up was judged out of scope alongside the
scale-up to ViT-B/L; H4 remains untested. Deferred to future work.

\section{Discussion}
\label{sec:discussion}

\paragraph{H1 (accuracy parity) --- not supported at this schedule.} None of
the three shared-projection variants comes within the \(\lesssim\)0.5 pt
target. In relative terms, Q-K=V's 2.71 pt drop (3.6\% relative) is close to
the 3.1\% perplexity cost reported for the analogous language-modeling
variant in \citet{kayyam2026three}, and QVV(3)'s 1.83 pt drop (2.5\%
relative) is similarly in that range --- but \citet{gao2026keyless} report
QVV(3) \emph{matching or exceeding} standard attention on language
perplexity, which does not reproduce here for ImageNet top-1. Q=K=V\(^+\)'s
7.62 pt drop (10.2\% relative) clearly exceeds the LM-scale gap. We cannot
rule out that a full 300-epoch DeiT schedule (rather than the 100-epoch
compressed one used here) would narrow these gaps, since shared-projection
variants may simply need more optimization steps to reach the same optimum.

\paragraph{H2 (overfitting reduction) --- untestable at this schedule, not
confirmed or refuted.} Every condition, including the QKV baseline, shows
\(|\Delta\text{Overfit}| \le 0.0014\) nats (Table~\ref{tab:imagenet1k}):
validation loss is still flat-to-falling at epoch 100 for all four models.
The DeiT-style augmentation stack (RandAugment, Mixup, CutMix, stochastic
depth, label smoothing) evidently suppresses overfitting for ViT-S/16 on
ImageNet-1k at 100 epochs regardless of attention variant, unlike the
36-layer language-modeling setting in \citet{gao2026keyless} where standard
attention showed a substantially larger gap (0.8404 vs.\ 0.6866 for QVV(3))
than any baseline here. Since the baseline itself doesn't overfit, there is
no gap for shared-projection attention to reduce, and this experiment cannot
distinguish a real absence of the effect from an effect masked by
regularization; a lower-augmentation or longer-schedule run would be needed
to test H2 properly.

\paragraph{H3 (symmetric vs.\ asymmetric penalty) --- mixed; collapse, not
symmetry, looks like the driver.} The fully symmetric Q=K=V\(^+\) trails the
asymmetric variants by 4.91 pts (vs.\ Q-K=V) and 5.79 pts (vs.\ QVV(3)) ---
a substantial penalty for collapsing to a symmetric attention map, running
counter to H3's expectation that spatial (non-causal) attention would make
symmetry comparatively cheap. The follow-up run (Section~\ref{sec:followup})
now supplies the second symmetric variant the original design called for:
Q=K-V\(^+\), which shares \(Q\) and \(K\) but keeps \(V\) separate, matching
Q-K=V's parameter count and cache target exactly. Its gap to Q-K=V is only
1.27 pts --- an order of magnitude smaller than Q=K=V\(^+\)'s 4.91-pt gap to
the same reference --- despite both Q=K-V\(^+\) and Q=K=V\(^+\) producing a
symmetric attention map. That points at collapsing \(V\) into the shared
projection, not attention-map symmetry per se, as the larger source of the
penalty, which is closer to H3's original motivation (symmetry itself being
comparatively cheap for spatial attention) than the verdict the single
Q=K=V\(^+\) data point supported alone. This reading carries the caveat
that Q=K-V\(^+\) trained at a different effective batch size than Q-K=V
(Section~\ref{sec:followup}), so it is suggestive rather than a controlled
confirmation, and a same-batch-size rerun is needed before treating H3 as
resolved either way.

\paragraph{H4 (detection transfer) --- untested.} No detection experiment
was run this iteration (Section~\ref{sec:results}); H4 remains open.

\section{Limitations}
\label{sec:limitations}
This iteration covers a single architecture (ViT-S/16), a single seed per
condition, and a single compressed schedule (100 epochs); none of these are
varied, so we cannot report variance across seeds or confirm that findings
hold at ViT-B/16 or ViT-L/16 scale, or under the full 300-epoch DeiT recipe.
The ImageNet-100 pilot phase and the ViT-B/16 confirmatory run described in
Section~\ref{sec:setup} were both skipped, as was the downstream
detection-transfer experiment (H4). All five originally-planned conditions
have now been trained, but not under matched settings: the follow-up run
adding Q=K-V\(^+\) (Section~\ref{sec:followup}) used effective batch 256 (2-GPU
DDP, linearly-scaled LR) rather than the batch 128 used for the other four
conditions, so its H3 comparison against Q-K=V is suggestive, not a
controlled replacement for a same-batch rerun. The same follow-up batch
added an unplanned sixth condition, Q=K=V with the positional correction
removed, whose large accuracy swing relative to Q=K=V\(^+\) is confounded
with that same batch-size difference and cannot yet be attributed to the
correction itself (Section~\ref{sec:followup}). As with \citet{kayyam2026three},
the compressed 100-epoch schedule (rather than the full 300-epoch DeiT
recipe) means absolute accuracy numbers should not be compared directly to
published SOTA ViT results trained with the full recipe or with
ImageNet-21k pretraining; the same compression is also the likely cause of
the null H2 result above, since it may not train long enough, or with light
enough effective regularization net of augmentation, for any condition to
visibly overfit.

\section{Conclusion}
\label{sec:conclusion}
We tested whether projection-sharing attention --- validated at
small-to-medium scale in \citet{kayyam2026three} and reframed and
generalized via depth-\(m\) factorization in \citet{gao2026keyless} ---
generalizes to ImageNet-1k, using ViT-S/16 as a first, deliberately narrowed
scale point. Q-K=V and QVV(3) recover most but not all of the QKV baseline's
top-1 accuracy (2.71 and 1.83 pt gaps respectively), while the fully
symmetric Q=K=V\(^+\) lags well behind (7.62 pt gap) --- a larger, not
smaller, symmetric-attention penalty than in the causal language-modeling
setting, running against H3's motivating expectation. A follow-up run
(Section~\ref{sec:followup}) filled in the previously-skipped Q=K-V\(^+\)
condition and found its gap to the matched-capacity asymmetric Q-K=V is only
1.3 pts, an order of magnitude smaller than Q=K=V\(^+\)'s --- suggesting the
earlier verdict conflated symmetry with the separate cost of collapsing
\(V\) into the shared projection, though this reading is not yet
confirmed under matched batch size. The same follow-up unexpectedly found
that removing the 2D positional correction from Q=K=V raises its top-1 by
about 4 pts, a result we can currently only report, not explain, since it is
confounded with that run's larger batch size. The
overfitting-reduction effect central to \citet{gao2026keyless}'s account of
\emph{why} shared-projection attention works did not appear for any
condition here, baseline included, because the DeiT augmentation recipe
suppressed overfitting broadly at this schedule; this leaves H2 open rather
than refuted. The ViT-B/16 confirmatory run, ImageNet-100 pilot,
detection-transfer evaluation, and a batch-size-matched rerun of the
follow-up conditions all remain as future work.

\bibliographystyle{plainnat}
\begin{thebibliography}{99}

\bibitem[Ainslie et al.(2023)]{ainslie2023gqa}
J.~Ainslie, J.~Lee-Thorp, M.~de~Jong, Y.~Zemlyanskiy, F.~Lebr\'on, and S.~Sanghai.
\newblock GQA: Training generalized multi-query transformer models from multi-head checkpoints.
\newblock In \emph{EMNLP}, 2023.

\bibitem[Borji(2023)]{borji2023kv}
A.~Borji.
\newblock Key-value transformer.
\newblock \emph{arXiv preprint arXiv:2305.19129}, 2023.

\bibitem[DeepSeek-AI(2024)]{deepseek2024mla}
DeepSeek-AI.
\newblock DeepSeek-V2: A strong, economical, and efficient mixture-of-experts language model.
\newblock \emph{arXiv preprint arXiv:2405.04434}, 2024.

\bibitem[Dosovitskiy et al.(2021)]{dosovitskiy2021vit}
A.~Dosovitskiy, L.~Beyer, A.~Kolesnikov, D.~Weissenborn, X.~Zhai, T.~Unterthiner, M.~Dehghani, M.~Minderer, G.~Heigold, S.~Gelly, J.~Uszkoreit, and N.~Houlsby.
\newblock An image is worth 16x16 words: Transformers for image recognition at scale.
\newblock In \emph{ICLR}, 2021.

\bibitem[Gao and Xu(2026)]{gao2026keyless}
X.~Gao and X.~Xu.
\newblock Keyless attention: Value-space routing and value-only caching for efficient transformers.
\newblock \emph{arXiv preprint arXiv:2606.21848}, 2026.

\bibitem[Kayyam et al.(2026)]{kayyam2026three}
A.~Kayyam, A.~Madan~Gopal, and M.~A.~Lewis.
\newblock Do transformers need three projections? Systematic study of QKV variants.
\newblock In \emph{Forty-third International Conference on Machine Learning (ICML)}, 2026.

\bibitem[Lee et al.(2021)]{lee2021rethinking}
J.~Lee et al.
\newblock Rethinking query, key, and value embedding in vision transformer under tiny model constraints.
\newblock \emph{arXiv preprint arXiv:2111.10017}, 2021.

\bibitem[Liu et al.(2021)]{liu2021swin}
Z.~Liu, Y.~Lin, Y.~Cao, H.~Hu, Y.~Wei, Z.~Zhang, S.~Lin, and B.~Guo.
\newblock Swin transformer: Hierarchical vision transformer using shifted windows.
\newblock In \emph{ICCV}, 2021.

\bibitem[MASA(2025)]{masa2025}
Anonymous.
\newblock Share your attention: Transformer weight sharing via matrix-based dictionary learning.
\newblock \emph{arXiv preprint arXiv:2508.04581}, 2025.

\bibitem[Shazeer(2019)]{shazeer2019mqa}
N.~Shazeer.
\newblock Fast transformer decoding: One write-head is all you need.
\newblock \emph{arXiv preprint arXiv:1911.02150}, 2019.

\bibitem[Touvron et al.(2021)]{touvron2021deit}
H.~Touvron, M.~Cord, M.~Douze, F.~Massa, A.~Sablayrolles, and H.~J\'egou.
\newblock Training data-efficient image transformers \& distillation through attention.
\newblock In \emph{ICML}, 2021.

\bibitem[Vaswani et al.(2017)]{vaswani2017attention}
A.~Vaswani, N.~Shazeer, N.~Parmar, J.~Uszkoreit, L.~Jones, A.~N.~Gomez, L.~Kaiser, and I.~Polosukhin.
\newblock Attention is all you need.
\newblock In \emph{NeurIPS}, 2017.

\end{thebibliography}

\end{document}